\section{Local comparison in families}\label{sec:local}

The common category $Q_S$ of Section~\ref{sec:walg} receives the two
semi-infinite pairings, one for $G$ and one for $\check G$.  We now compare
these pairings over the entire space of marked points, including its
collision diagonals.  The point calculation supplies a gap in low
cohomological degrees.  That gap determines the equivalence of hearts
across a collision; derived realization then supplies the equivalence of
DG categories.  A further argument with translation of the marked points
will make the comparison compatible with insertion in both slots.

Throughout this section $a$ is irrational and $d$ is its dual shifted
level.  The spherical category $C_S$, the test category $d_S$, and their
pairing $F_S$ are those of Section~\ref{sec:loop}, with
$d_S=\mathcal W_{-s}^-(S)$ as in
\eqref{loop:whittaker-coinvariants}.  In particular, $d_S$ is a category of tests,
whereas $d$ without a subscript is a level.  Checked symbols refer to the
dual group.  We work on smooth affine coordinate charts for tuples of
points.  Extra auxiliary coordinates are allowed.  When adding a
collision diagonal, we retain the complements of any other diagonals
already deleted; thus the new diagonal is a smooth divisor and its
complement is principal affine locally.

\subsection{The heart and its derived realization}

We use the $t$-structure on $C_S$ whose heart has the perverse
$D$-module normalization in the parameter $S$.  In the underlying
quasi-coherent description $b$, its test includes the shift $[\dim S]$.
This convention is relevant to both insertion and the open--closed
arguments below.

\begin{proposition}\label{local:heart}
On the coordinate charts just described, the $t$-structure on $C_S$ is
separated and compatible with filtered colimits.  Its heart is a
Grothendieck abelian category, and there is a continuous equivalence
\begin{equation}\label{eq:L8}
                     D(C_S^\heartsuit)\simeq C_S.
\end{equation}
The functors $j^!,j_*,i_*$ are $t$-exact when $j$ is a principal open and
$i$ is a smooth closed collision diagonal.  For a smooth divisor
$i$, the functor $i^!$ has amplitude $[0,1]$.  These assertions also hold
for products of the group categories, with separate subsets of the
marks assigned to the different factors.
\end{proposition}

\begin{proof}
We first identify the connection that must be added to the
Harish--Chandra description of Section~\ref{sec:loop}.  On an affine
smooth $S$, forgetting the connection has the continuous left adjoint
of $D$-induction.  Its projection formula allows this adjunction to be
tensored with a crystal of categories.  Forgetting is conservative:
if the underlying object of $M$ vanishes, the action of the
$D$-induced structure sheaf on $M$ vanishes, and this induced object
generates $D(S)$.  The adjunction is therefore monadic.

The connection monad is the differential-operator monad with crystal
transport.  Filter it by the order of distributions on the formal
diagonal of $S$.  At every finite order the transport is the one over
that infinitesimal diagonal; the associated graded is tensoring with
$\operatorname{Sym}(T_S)$.  Combine this filtration with the
Harish--Chandra induction monad of Proposition~\ref{loop:hc-renormalization}.
In coordinates, the combined induction adds the Laurent-current
action and the derivatives on $S$, the latter taken at fixed curve
coordinate $z$.  Regular currents are stable under this derivative.
For a distribution of fixed order, the order of a vanishing current
can decrease only by a bounded amount.  Thus the formula respects
algebraic arc equivariance and smoothness.  Each associated graded
piece for distribution order is the induction in \eqref{eq:L2}
tensored with a vector bundle.  It is $t$-exact.  The increasing
filtration and passage to Ind extend this conclusion from bounded
below arc representations to the renormalized category in both
$t$-directions.

Let $\mathcal H_S$ denote the ordinary abelian category of smooth
connection Harish--Chandra modules obtained in this way, with the
specified parameter shift.  The two consecutive forgetful functors
are conservative and continuous, so their monadic descriptions give
a $t$-exact comparison
\[
                       u_C:C_S\longrightarrow D(\mathcal H_S)
\]
which is an equivalence on bounded below objects.  The heart
$\mathcal H_S$ is Grothendieck.  It is generated by the modules induced,
including their connection, from the arc representations that are
locally free over the base.  The same induction adjunctions give
bounded compact generators of $C_S$.  Forgetting further to underlying
$D$-modules, and removing the parameter shift when using left crystal
normalization, commutes with the parameter operations.

We verify conservativity also without a boundedness condition.
At a point the category $C$ is DG semisimple, with basis the Weyl
modules $M_\eta$ for actual dominant weights of $G$.  Here is the
irrational-level argument.  The Sugawara operator $L_0$ on $M_\eta$
has bottom energy equal to the finite Casimir divided by $2a$, and
its spectrum consists of this number plus nonnegative integers.
A submodule with a higher bottom energy has, at its bottom, a finite
$G$-type killed by the positive currents.  Its energy difference
would be a nonzero rational number divided by $2a$ and a positive
integer.  Irrationality excludes this.  A submodule reaching the
bottom contains the generating irreducible finite $G$-module.
Hence each $M_\eta$ is simple.  Extensions split by energy separation;
at a common bottom energy, a splitting of finite $G$-modules lifts
the bottom vectors, which the positive currents kill, and induction
splits the extension.  The statements use the algebraic energy
decomposition from PBW, and generalized eigenspaces for an extension.
The unipotent filtration of a finite arc representation now shows
that its induction is a sum of these Weyl modules.  These inductions
generate the heart.  Bounded comparison and generation by bounded
compacts give the asserted DG semisimplicity.

On a cell of fixed coincidence pattern the coordinate description
of $C_S$ is the tensor product of the corresponding point categories
with $D(S)$.  This description respects the crystal: infinitesimal
translation of centered coordinates is inner, through Sugawara
translation, since $a\ne0$.  Between nilpotent lifts its exponential
is a finite expression, acting on modules and their mappings.
Consequently $u_C$ detects objects on each cell.  The ! restrictions
to the cells are jointly conservative by the open--closed triangles,
so $u_C$ detects objects on the chart.  This proves separation.
The claimed exactness and amplitude of the parameter operations
can therefore be checked in ordinary $D$-modules.  In particular,
principal-open star extension is exact; the localization triangle
then gives the amplitude of $i^!$.  Compatibility with filtered
colimits follows from the monadic formulas.

It remains to justify unbounded realization.  We describe the
construction, since bounded comparison alone would not suffice.
An exact, sum-preserving functor from a Grothendieck heart to the
heart of a separated DG category with filtered-colimit-compatible
$t$-structure can be realized on complexes as follows.  For a
bounded above complex, shift its upper bound to degree zero and
use the additive Dold--Kan construction, followed by geometric
realization in the target.  Equivalently, use the filtered colimit
of the finite totalizations of its lower stupid truncations.
Changing the upper bound gives the same construction by the
normalized-chain suspension identification.  For an arbitrary
complex, take the filtered colimit over its good upper truncations.

The inclusion of the heart has discrete mapping spaces: negative
cohomology of a mapping complex between heart objects vanishes.
Thus the diagrams and their coherences used here are actual diagrams
in the target.  Short exact sequences of heart objects are cofiber
sequences there.  Finite totalizations compute the usual cohomology;
filtered-colimit compatibility gives the same assertion for the
successive infinite totalizations.  In particular, quasi-isomorphisms
are inverted.  The cone calculation is compatible with changing the
truncation bounds, since the resulting comparison is an isomorphism
on cohomology and the $t$-structure is separated.  This constructs an
exact continuous DG realization functor.

Apply this construction to the inverse of $u_C$ on hearts.  Its
composition with $u_C$ is the identity by the same realization formula
in the ordinary derived category.  The opposite composition is the
identity on bounded objects by bounded comparison, and hence on $C_S$
by its bounded compact generators and continuity.  This proves
\eqref{eq:L8}.  The construction is linear: tensoring with a complex
of vector spaces is computed by the same finite operations, sums,
and realizations.
\end{proof}

We will also use the corresponding statement about natural
transformations.  Suppose two continuous functors out of
$D(\mathcal H_S)$ have images of heart objects in a common subcategory
whose mapping spaces are discrete.  A natural transformation on the
heart extends uniquely by the realization formulas above.  Indeed,
apply it term by term to Dold--Kan diagrams and their truncation
colimits.  An existing transformation must equal this extension by
naturality with the colimit cocones.  This argument identifies the
mapping spaces as well as the sets of transformations.  It applies
separately in each variable of a multilinear functor, including
auxiliary scheme $D$-module variables.

\subsection{Insertion of marked points}

Let $T\subset V$ be two assignments of marks over the same base $S$.
All collisions remain allowed.  Subscripts $T$ and $V$ indicate the
corresponding categories, and $T=\varnothing$ means the unit input
$D(S)$.

\begin{lemma}\label{local:insertion}
Insertion of the unit defines $t$-exact functors
$C_T\to C_V$ and $d_T\to d_V$.  They have continuous $D(S)$-linear
right adjoints, commute with parameter operations and disjoint
products, and use the vacuum when $T=\varnothing$.
On arc-induced modules the first functor is
\begin{equation}\label{eq:L9}
 \operatorname{Ind}_{j_T}^{\ell_T}W
       \longmapsto
 \operatorname{Ind}_{j_V}^{\ell_V}(W|_{H_V}).
\end{equation}
Here restriction is along the restriction map on regular arcs.
\end{lemma}

\begin{proof}
The Lie-algebra formula factors through currents over $V$ whose
poles are allowed only at $T$.  Modulo regular currents, this
intermediate algebra has the same principal parts as the algebra
for $T$.  Inducing further gives \eqref{eq:L9}, the usual insertion
formula for localization; compare \cite{GLCII}.  The formula is
horizontal for the fixed-$z$ connection of
Section~\ref{sec:loop}.  It preserves compact generators.  Its
right adjoint is therefore continuous, and the quasi-coherent-linear
adjunction descends to the crystal.  Compact generation and the
mapping base-change statement show that the adjoint commutes with
base change as well.

We check exactness even when points collide.  Add one new point $x$
and write $w=z-x$.  If $f$ is the old divisor equation, the quotient
of Laurent currents by the intermediate algebra is exhausted by
principal parts in $w$ after inverting $f$.  At finite order the
underlying module is built from
$\mathcal O/(w^n)[f^{-1}]$, where $\mathcal O$ is the completion
for all the marks.  The quotient $\mathcal O/(w^n)$ is the finite
jet ring at $x$.  Inverting $f$ in it is equivalent to inverting
$f(x)$ in the base: their difference is nilpotent modulo $w^n$.
These modules, and hence the quotient of current algebras, are flat
over $S$.

For completeness, flatness suffices for PBW induction even if the
quotient is not projective.  Modulo the common regular lattice, the
sequence of principal-part modules is pure exact with flat quotient.
In characteristic zero the symmetric-power filtration for this
sequence is obtained from the tensor-power filtration by the
symmetric idempotent.  The associated graded induction therefore
adds a flat symmetric-algebra factor.  Alternatively, resolve
induction by the common-lattice free bar and filter by PBW degree.
For a heart input its associated graded is precisely this
symmetric-algebra bar, concentrated in heart degree zero.
Induction on the filtration and exhaustion prove exactness of the
filtered induction.  Proposition~\ref{local:heart} transfers this
calculation to $C_S$.  Iterating adds any finite set of marks;
for the empty input the assertion is ordinary PBW.

For $d_T\to d_V$, use the opposite-character Whittaker model
\eqref{loop:whittaker-coinvariants}.  On each pole bound the functor is the closed push
of the locus with no added poles.  Its continuous extension is
$t$-exact, and its right adjoint is ! restriction.  Both preserve
the Whittaker conditions: on the relevant strata the new distinct
points have zero pole contribution, with the same $q_D$-pull and
exponential; negative defect has zero restriction.  Equivalently,
one may use the projection characterization of these conditions
after \eqref{eq:W1}.  Empty insertion is the zero-pole vacuum
locus, whose Whittaker category is $D(S)$ and whose restriction
adjoint is continuous.  By ! base change these closed pushes also
agree with insertion of the Grassmannian kernels in the
$\mathbf d$ description of \eqref{eq:loop-whittaker-functional}.  All the formulas
respect disjoint products and parameter operations.
\end{proof}

\subsection{The point calculation and its cohomological gap}

At one point, write $M_\eta$ for the Weyl module with dominant weight
$\eta$, and $D_\lambda$ for the opposite-character standard in
\eqref{eq:W3}, labeled by a dominant coweight $\lambda$.  The labels
are those of the actual root datum; no passage to a simply connected
or adjoint replacement is made.

\begin{lemma}\label{local:point-gap}
The point pairing satisfies
\begin{equation}\label{eq:L10}
 R\operatorname{Hom}_{Q}
 \bigl(F(M_\eta,D_\lambda),F(M_{\eta'},D_{\lambda'})\bigr)
 \simeq
 \begin{cases}
 H^\bullet_{\mathrm{dR}}(G),&
                    (\eta,\lambda)=(\eta',\lambda'),\\
 0,&\text{otherwise}.
 \end{cases}
\end{equation}
The map induced from the input categories identifies cohomology
in degrees at most two.  Under the common $W$-module comparison,
$F(M_\eta,D_\lambda)$ is the irrational-level simple reduction
$T^\kappa_{\eta,\lambda}$, tensored with a one-dimensional vector
space.  Its essential image agrees with the image for the dual
pairing with the two labels exchanged.
\end{lemma}

\begin{proof}
Let $q$ be the balanced-tensor projection in \eqref{eq:L3}.
The monad $q^Rq$ is obtained by applying the right adjoint of
multiplication in the spherical category $\mathcal S$ to its unit,
then acting on the two inputs.  At a finite quotient $H_f$ of the
arc group by a pro-unipotent kernel this right adjoint uses
\[
              \Delta_{BH_f,*}\omega_{BH_f}[-2\dim H_f].
\]
Pulling to the two trivial $H_f$-bundles gives
$H^\bullet_{\mathrm{dR}}(H_f)=H^\bullet_{\mathrm{dR}}(G)$.
On $C$, the spherical action supported at the unit orbit tensors
the underlying Lie module by its ! fiber at the base point of
$BH$.  The semisimplicity proved in Proposition~\ref{local:heart}
therefore computes this action on distinct Weyl labels.

For the test slot, put
$a_\lambda=\langle2\rho_{\mathrm{rt}},\lambda\rangle$ and let
$\mathbf k_g$ be the kernel at the section $g=t^\lambda$.
Equations~\eqref{eq:W2} and \eqref{eq:W3}, including integration
in dimension $\delta_N+a_\lambda$, give, after retaining the
constant one-dimensional line,
\[
                   \mathbf d(\mathbf k_g)=D_\lambda[-a_\lambda].
\]
Unit-orbit convolution on this section also contributes the fiber
at the base point of $BH$: the section trivializes its arc torsor.
Combining the two actions proves \eqref{eq:L10}.  The unit of the
adjunction is the constants map.  Since $G$ is semisimple,
$H^1_{\mathrm{dR}}(G)=H^2_{\mathrm{dR}}(G)=0$, and the positive
cohomology starts in degree at least three.

We now use the irrational-level reduction theorem of
Arakawa--Frenkel~\cite[Theorems~2.1--2.3]{ArakawaFrenkel}.
Principal reduction of the
Weyl module labeled by $\eta$, with Drinfeld--Sokolov twist by the
dominant coweight $\lambda$, is concentrated in degree zero and
is the simple highest module $T^\kappa_{\eta,\lambda}$.  Under
Feigin--Frenkel duality it identifies with
$T^{\kappa^\vee}_{\lambda,\eta}$.  In our modes the character on
a simple-root current is at mode
$-1-\langle\alpha,\lambda\rangle$, the lattice is the regular
lattice, and the $W$-action is the spectral-flow action.  Simultaneous
changes of character sign conjugate the $W$-fields and leave the
statement unchanged.

To match these conventions, apply \eqref{eq:L1} to $\mathbf k_g$.
Spectral flow changes a simple-root mode $m$ to
$m-\langle\alpha,\lambda\rangle$.  Its conjugated Fock lattice
differs from the regular lattice by the determinant line and the
shift $[-a_\lambda]$ in \eqref{eq:L6}.  This shift cancels the
shift in the displayed formula for $\mathbf d(\mathbf k_g)$.
The resulting ordinary compared module is exactly the reduction
in the preceding paragraph.  These objects are compact by the
adjunctions of \eqref{eq:L3}; full faithfulness of the ordinary
$W$ comparison on compacts promotes this calculation to $Q$.

We will need its generating vector at a collision.  It is the
highest Weyl vector times the semi-infinite vacuum for the specified
lattice.  The character and highest-vector formulas make it closed.
In the DS calculation, coinvariants for the strictly negative
$N$-modes evaluate the PBW-free negative unipotent factor at the
character, so the finite-type generating vector survives.  The
remaining nonnegative cochain calculation has no degree-zero
boundaries.  Compare also the highest-weight and grading calculations in
\cite[Sections~4.1 and~4.3]{ArakawaFrenkel}.  Thus its class is nonzero and
generates the reduction.  No rigidification of its one-dimensional
line is needed.
\end{proof}

On a cell of fixed coincidence pattern, the same calculation works
with parameters.  Indeed $Q_S$ is locally the point category tensored
with $D(S)$.  Its horizontal identification is supplied by $W$
Virasoro translation: between formal lifts, exponentiation of the
translation generator in a nilpotent ideal identifies centered
coordinates.  This is an infinitesimal crystal calculation, and
requires no exponentiation on a punctured parameter space.
The image of a pair of labels is consequently its fixed compact
point object tensored by a line with connection.  To see that only
a line can occur, use mapping base change and \eqref{eq:L10}:
negative self-Ext vanishes and the degree-zero endomorphisms are
scalars.  The vacuum line is canonically the same for the two
pairings.  In \eqref{eq:L6} its generating vector is the Weyl
vacuum times the $h$-Fock vacuum, horizontal for the absolute
currents; the vacuum presentation in \eqref{eq:L7} preserves it
over families and their collision diagonals.

The same description computes right adjoints.  On coordinate
pieces of a cell the pairing is tensoring with fixed compact point
objects and invertible line connections.  Therefore the cone of
its unit on a heart input lies in degrees at least three, and
every input is a retract of the result of applying the resulting
right-adjoint monad.  Products of points give the same assertions.
Insertions simply prescribe that some labels be the vacuum.

\subsection{Extending the equivalence across collisions}

Insert the vacuum in the second slot of $F_S$ and in the first
slot of $\check F_S$, and write
\[
 F_0:C_S\longrightarrow Q_S,\qquad
 F_1:\check d_S\longrightarrow Q_S.
\]
Their right adjoints $R_0,R_1$ are continuous and $D(S)$-linear,
by Proposition~\ref{loop:pairing-adjoints} and
Lemma~\ref{local:insertion}.  Corollary~\ref{loop:parameter-operations}
identifies these adjunctions under restriction to the smooth parameter
charts.  In particular they commute with $i_*,i^!,j_*,j^!$ for the
closed collision divisor and its principal open complement, including
the units and counits.  The same assertion applies to closed graph
maps in auxiliary parameter factors.

\begin{proposition}\label{local:one-slot}
There is a canonical continuous $t$-exact equivalence
\begin{equation}\label{eq:L11}
 \alpha_S:C_S\xrightarrow{\ \sim\ }\check d_S,
                 \qquad F_1\alpha_S\simeq F_0.
\end{equation}
It is $D(S)$-linear and compatible with the parameter operations,
disjoint products, and coordinate changes.  Exchanging the two
groups gives $\check\alpha_S:\check C_S\simeq d_S$.
The same statements apply to the simultaneous opposite levels.
\end{proposition}

\begin{proof}
The proof has three stages.  We propagate two low-degree properties
across a collision divisor, use them to construct the equivalence of
hearts, and then realize that equivalence and its compatibilities.
The final passage to DG categories will also use the cell retract
property established after Lemma~\ref{local:point-gap}.

On a cell, that lemma gives the following properties for heart inputs:
\begin{enumerate}
\item \emph{Unit gap.}  The cone of $M\to R_rF_rM$ lies in degrees
      $\ge3$, for $r=0,1$ in its respective source category.
\item \emph{Cross-pairing gap.}  For $U=R_1F_0$, the object $UM$ lies in degrees $\ge0$,
      has $H^1=H^2=0$, and $H^0U$ is an exact equivalence of
      hearts.  The adjunction map
      $F_1H^0UM\to F_0M$ is an isomorphism.
\end{enumerate}
The analogous assertions hold with the subscripts reversed.
The invertible base lines in the cell calculation are included
in these equivalences.

\smallskip
\noindent\emph{Propagating the unit gap.}
We prove these properties on a larger piece of the arrangement by
adding a smooth divisor $i$, with principal complement $j$.
Assume them on both smaller pieces.  For $M$ in the heart put
$B=j^!M$ and $E=i^!M$.  Then $B$ is in the heart and $E$ has
amplitude $[0,1]$, by Proposition~\ref{local:heart} and
Proposition~\ref{whittaker:t-structure} on the two respective sides.
Apply the cone of the unit to
\[
                  i_*E\longrightarrow M\longrightarrow j_*B.
\]
The exact pushes and the known bound on the two cohomology terms
of $E$ give the unit gap on the larger piece.  It implies, by
adjunction, that each $F_r$ identifies degree-zero maps between heart inputs with
those between their images, whose negative Hom also vanishes.

\smallskip
\noindent\emph{The boundary map and the equivalence of hearts.}
On the closed piece, the induction hypothesis gives $UE$ in
degrees $\ge0$, with $H^0UE,H^1UE$ obtained from $H^0E,H^1E$
by the heart equivalence, and $H^2UE=0$.  Moreover,
\[
                        F_1\tau_{\le1}UE\simeq F_0E.
\]
Indeed this holds on the two heart terms; the gap makes their
truncations a triangle, so it holds on $E$.

Apply $U$ to the localization triangle.  Its only possible
obstruction to the cross-pairing gap is $H^1UM$.  In the original heart,
the boundary gives an epimorphism
\[
                    j_*B\longrightarrow i_*H^1E.
\]
The corresponding boundary after applying $U$ is
\[
                j_*H^0UB\longrightarrow i_*H^1UE.
\]
To identify these maps, let $D$ be a heart object on the closed
piece of the $C$ family.  Here and below $U$ on either smaller
piece denotes its corresponding adjoint composite.  There are
natural identifications of degree-zero maps
\[
\begin{aligned}
 \operatorname{Hom}_{C_S}(j_*B,i_*D)
 &\simeq \operatorname{Hom}_{Q_S}(F_0j_*B,F_0i_*D)\\
 &\simeq \operatorname{Hom}_{Q_S}
       (F_1j_*H^0UB,F_1i_*H^0UD)\\
 &\simeq \operatorname{Hom}_{\check d_S}
       (j_*H^0UB,i_*H^0UD).
\end{aligned}
\]
The outer identifications use the unit gaps just proved on the
larger piece; the middle one uses the smaller-piece comparisons
and their compatibility with the exact pushes.  Apply this chain
with $D=H^1E$.  The comparison
$F_1\tau_{\le1}UE\simeq F_0E$, followed by the truncation map
$E[1]\to H^1E$, identifies the two boundary maps.  The gaps
ensure that the maps to the indicated heart objects are unique.
The cokernel of the second boundary is supported on $i$, because
$j^!$ is exact.  By the equivalence of hearts on the closed piece,
write it as $i_*H^0UD$ for a closed-heart object $D$.  The quotient
map from $i_*H^1UE$ then corresponds to a map
$i_*H^1E\to i_*D$.  Naturality of the displayed identifications
makes its composite with the first boundary zero.  That boundary
is an epimorphism, so this map, and hence the quotient map, is zero.
The cokernel therefore vanishes.  Hence
$H^1UM=0$; the other degree bounds and $H^2UM=0$ follow from
the same triangle.

Finally, test the adjunction map $F_1H^0UM\to F_0M$ by !
restriction to both pieces.  On the open it is the known map.
On the closed piece the gap and the amplitude of $i^!$ give
\[
                         i^!H^0UM=\tau_{\le1}UE,
\]
so it is the preceding isomorphism for $E$.  Conservativity of
the two restrictions proves the assertion.  The long exact
sequence and the gap show that $H^0U$ is exact on the heart.
Repeating with the two sides reversed gives inverse heart
functors, because the unit gap identifies their degree-zero
maps.  This completes the induction through the arrangement.

\smallskip
\noindent\emph{Realization and equivalence of DG categories.}
Realize $H^0U$ using \eqref{eq:L8}.  The heart formula commutes
with sums and filtered colimits, so its realization $\alpha_S$
is continuous and $t$-exact.  The comparison of functors in
\eqref{eq:L11} extends by the natural-transformation statement
after Proposition~\ref{local:heart}: the relevant images of
heart objects have discrete mapping spaces.  This also proves
commutation with the exact operations $j^!,j_*,i_*$, first on
hearts and then on their realizations.

We verify that $\alpha_S$ is an equivalence without presupposing
a derived-realization theorem for $\check d_S$.  The mate
\[
                            \alpha_SR_0\longrightarrow R_1
\]
is an isomorphism on the cells, by their explicit description.
It remains so after successive $i_*$ and $j_*$, by their linear
adjunctions and the comparisons just established.  Open--closed
triangles show that it is an isomorphism on all objects.
Adjunction and \eqref{eq:L11} now give full faithfulness on maps
into objects in the image of $R_0$.  On every cell, every input
is a retract of such an object by the cell monad calculation.
Applying the exact pushes and the localization triangles yields
full faithfulness everywhere.  The corresponding retracts on
the other side generate $\check d_S$ under colimits and finite
extensions.  Thus the continuous fully faithful functor is also
essentially surjective.

\smallskip
\noindent\emph{Linearity and coherence.}
Carry out all the
constructions with a list of auxiliary scheme factors.  Exterior
products of heart objects are heart objects in the ordinary
module description; for the Whittaker side this follows on the
pole bounds and then by colimits.  The same gaps therefore give
the canonical comparisons for exterior products.  The comparisons
also commute with closed graph or diagonal pushes in these
auxiliary factors, by the exact heart constructions.  Their
coherences are unique there, since the relevant mapping spaces
are discrete, and realization preserves these coherences.
Taking mates gives the coherent squares for ! pulls.  Exterior
product followed by diagonal ! pull computes the $D(S)$ action;
hence these squares give $D(S)$-linearity, including its unit
and iterated-action coherences.  The original $F_i,R_i$ commute
with the units and counits of these operations, so the comparison
$F_1\alpha_S\simeq F_0$ has the same linearity.

Disjoint products give the cell comparison used to construct
$\alpha_S$, and hence the same comparison everywhere.  To check
the $t$-structures even when each disjoint block has internal
collisions, characterize nonnegative objects by their !
restrictions to the arrangement pieces, using the localization
triangle and left exactness of ! restriction and the pushes.
The product equivalence and its inverse preserve this criterion.
Uniqueness on hearts and realization give its coherence.
The same criterion makes the $t$-structures local in the curve,
including those defined by global pole models for $d_S$.
Consequently the constructions agree on étale coordinate charts
and glue.  This proves all the assertions.
\end{proof}

\subsection{Comparison with both slots inserted}

The preceding equivalence will be used inside integrated pairings.
For this purpose it must respect insertion simultaneously in both
slots.  Let $T,U$ be two colors of independent marks over a parameter
chart $S$, let $Z$ be extra neutral marks, and put
$V=T\cup U\cup Z$.  Collisions within and between all three sets
are allowed.  The common input is
$C_T\otimes_{D(S)}\check C_U$, with $D(S)$ used for an empty color.

\begin{proposition}\label{local:two-slot}
There is a canonical comparison $G_0\simeq G_1$, where
\begin{equation}\label{eq:L12}
 \begin{aligned}
 G_0(A,B)&=F_V\bigl(\operatorname{ins}A,
                   \operatorname{ins}(\check\alpha_U B)\bigr),\\
 G_1(A,B)&=\check F_V\bigl(\operatorname{ins}B,
                   \operatorname{ins}(\alpha_T A)\bigr).
 \end{aligned}
\end{equation}
It agrees with the disjoint comparison and respects $D(S)$-linearity,
coordinate changes, disjoint products, symmetries, and all diagonal
operations for merge and repetition of marks.  Vacua give the empty
inputs.  In particular, \eqref{eq:L11} is compatible with ordinary
insertion of marks.
\end{proposition}

\begin{proof}
First the collision argument constructs an autoequivalence of the
common input that intertwines $G_0$ and $G_1$.  Translation will
identify this autoequivalence with the identity.  We will then check
that the resulting comparison specializes to the prescribed vacuum
insertion map when a neutral point meets another mark.

Both functors have continuous linear right adjoints by
Proposition~\ref{loop:pairing-adjoints} and
Lemma~\ref{local:insertion}; their parameter adjunctions are those of
Corollary~\ref{loop:parameter-operations}.  To apply the proof of
Proposition~\ref{local:one-slot}, take both source families to be
$C_T\otimes_{D(S)}\check C_U$ and replace $F_0,F_1$ by $G_0,G_1$.
Proposition~\ref{local:heart} supplies the product KL $t$-structure,
its open--closed amplitudes, and derived realization.  On each cell
the point label calculation supplies the unit and cross-pairing gaps,
as well as the retracts of right-adjoint images used for the derived
equivalence.  This holds also at cross-color collisions: the point
calculation exchanges the matched labels, while the neutral
insertions designate vacua.  The divisor induction, realization,
and auxiliary-parameter argument therefore give a $t$-exact
$D(S)$-linear autoequivalence $\beta_S$, whose heart functor is
$H^0(G_1^R G_0)$, and an isomorphism
\begin{equation}\label{eq:L13}
                           G_1\beta_S\simeq G_0.
\end{equation}
On the open where different colors do not meet, neutral marks meet
neither color, and neutral marks are distinct, this autoequivalence
is canonically the identity.  Use disjoint factorization,
\eqref{eq:L11}, and the common horizontal $W$ vacuum.  This comparison
retains the line connections; it does not choose flat frames for
conformal-weight lines.  We show that the identity comparison extends
through the deleted crossings.

\smallskip
\noindent\emph{Translation of an entire color.}
We construct the comparison first on the universal affine-line chart
with independent coordinates for the marks and any auxiliary
parameters.  The parameter comparisons then pull it to the smooth
charts under consideration; curve-coordinate descent is checked
below.  If $T$ is nonempty, select one of its marks as origin for
the relative coordinates and write
\[
                           S=\mathbb A^1_y\times S_0,
\]
where $y$ translates all $T$-marks together.  In this step remove
only those collision loci not involving $T$, so their equations
are independent of $y$.  The locus of forbidden crossings involving
$T$ is the zero set of a polynomial $f$ monic in $y$.
The common input admits a $t$-compatible presentation
\[
           D(\mathbb A^1)\otimes\mathcal E(S_0).
\]
Here $\mathcal E(S_0)$ is the product KL category for the relative
positions and the other color; it has the heart and realization
properties of Proposition~\ref{local:heart}.

We justify the horizontal assertion in this presentation also at
internal collisions of $T$.  Simultaneous translation is implemented
by $L_{-1}$, the total smeared Sugawara operator.  Its bracket with
currents is the ordinary translation derivation.  Use total residue
over the marked divisor to define it.  The mode expansion and the
bracket identity hold over the parameter ring, as can be checked on
PBW generators before specialization.  In particular this operator
is horizontal in the remaining parameters at fixed curve coordinate
$z$.  One can see this directly by splitting into the regular
lattice and the isotropic Laurent complement of rational functions
with poles on the divisor and vanishing at infinity.  Infinitesimal
variation of the poles at fixed $z$ preserves both halves.  The
quadratic expression from their residue pairing, with regular modes
on the right, has the translation bracket on all the fibers and
over the parameter ring.  Normal-order sums act on smooth modules.
Subtracting this inner translation from the parameter derivative
gives the commuting derivative in the displayed product category.
The nilpotent exponential identifies the centered and absolute
crystal transports, including on formal collision diagonals.
The same constructions allow the fixed origin among the relative
coordinates of $S_0$.

\smallskip
\noindent\emph{Unique extension on constant objects.}
For $E\in\mathcal E(S_0)^\heartsuit$ set
\[
                        M=\omega_{\mathbb A^1}[-1]\boxtimes E.
\]
This is a heart object with no nonzero subobject or quotient
supported on $f=0$.  For subobjects this follows from the ordinary
module presentation: the polynomial $f$ is monic, so multiplication
by it is injective on the module underlying $M$.  For quotients,
a generator constant along $y$ cannot have nonzero image in an
$f$-torsion module.  If $f^n e=0$ and the commuting derivative
satisfies $\partial_y e=0$, differentiation
$n\deg_y(f)$ times gives a nonzero constant multiple of $e=0$.
The argument applies generator by generator, without coherence
or a uniform torsion exponent.  Vanishing on the open $f\ne0$
is exactly $f$-torsion in this module description.

The same condition holds for $\beta_S M$.  Indeed $\beta_S$ is
an exact equivalence, commutes with open restriction along with
its inverse, and hence preserves the supported subcategory.
The identity comparison over $f\ne0$ now extends uniquely to
$\beta_S M\simeq M$.  To prove this explicitly, embed both
objects in the star push of their common open restriction; the
embeddings are monomorphisms because neither object has a
supported subobject.  The quotient of either image by its
intersection with the other is supported on $f=0$, hence zero.
Their images are equal.  The same argument proves uniqueness of
the extension.

This comparison is natural in $E$.  Extend it by realization in
$\mathcal E(S_0)$, and then by $D(\mathbb A^1)$-linearity.  The latter
step is valid because the displayed product is a free
$D(\mathbb A^1)$ factor and $\omega_{\mathbb A^1}[-1]$ is its shifted
unit.  To retain $D(S_0)$-linearity, do the construction with any
list of auxiliary parameters.  Uniqueness on the external heart
objects makes it commute with exterior products and with exact
graph and diagonal pushes in those parameters.  Realization and
the free $D(\mathbb A^1)$ factor extend these identities.  Passing
to mates supplies the coherent !-pull and linearity squares,
as in Proposition~\ref{local:one-slot}.

This allows all crossings involving $T$.  Next translate $U$ to
allow its remaining crossings.  Finally translate one neutral
point at a time to allow crossings between neutral points; the
input then has a free $D(\mathbb A^1)$ factor with no color on it.
At each step the deleted equation is monic in the translating
coordinate, while the other remaining deletions lie in the base.
We have constructed an extension of $\beta_S\simeq\operatorname{Id}$
over every collision.

\smallskip
\noindent\emph{Uniqueness and changes of chart.}
Natural transformations after \eqref{eq:L13} are controlled on
hearts by Proposition~\ref{local:heart}.  An endomorphism of the
identity of a KL input is determined on the dense open.  Indeed,
its heart is epi-generated by connection-induced arc
representations locally free over the base.  PBW makes these
inducing generators torsion-free for localization to the dense
open.  A map on such a generator that vanishes there is zero;
naturality with the generating epimorphisms then detects the
transformation on every heart object.  The same argument with
auxiliary parameters and diagonal pushes proves coherent
uniqueness.  Thus the extensions just constructed are independent
of the order of translation, agree on smaller affine charts, and
glue.  Over $X$, use étale coordinates separating disjoint
clusters.  The locality in the curve proved in
Proposition~\ref{local:one-slot} identifies these constructions
with the affine-line calculation.  This gives \eqref{eq:L12},
including its symmetries and disjoint-product compatibility.

\smallskip
\noindent\emph{Diagonal operations involving a neutral point.}
The comparison now extends over the independent marking chart.  To
identify its pullback along a neutral-point diagonal with the
specified insertion comparison, we must still determine its value
on the vacuum.  The following two-point calculation does this.

Compatibility with a same-color merge follows from uniqueness:
on the resulting scheme use its dense open with separated
colors and neutrals, where the comparison is \eqref{eq:L11}.
For a neutral point merging into another mark the same uniqueness
reduces to the locus where all other resulting marks are separated.
These reductions are legitimate inside the comparison of functors:
closed diagonal push is fully faithful and $t$-exact, and the gap
and adjunction constructions commute with it.  Thus one can push
a resulting heart object to the original chart, compare there,
and apply the same heart and realization test on the diagonal.

Disjoint product reduces the remaining calculation to two points:
a neutral point $w$ merging with a neutral point or a labeled
point $x$.  Two neutral points use the common horizontal vacuum
already described.  Now let the label at $x$ have color $T$.
The first pairing uses its Weyl module and the inserted test vacuum;
the second uses the dual Weyl vacuum and the matched test label.
Write $\lambda$ for the coweight of that dual test label; the $U$
case exchanges the groups.  Work in a coordinate neighborhood and
fix $x$ by base change if convenient.  On the
Weyl side, insertion is \eqref{eq:L9} from evaluation at $x$.
On the flow side, use the Grassmannian section
$g=(z-x)^\lambda$ and the inserted Weyl vacuum.  Retain the shifts
and determinant line from Lemma~\ref{local:point-gap}.  The line
comparison in \eqref{eq:L11} pulls back from the label at $x$.

In the compact ordinary comparison \eqref{eq:L6}, both complexes
have the closed cycles given by their highest vectors and Fock
vacua.  The mode formulas \eqref{eq:L7} define them also at $w=x$.
On the flow side the change of lattice has a finite locally free
quotient at $x$, even when $w=x$, so its determinant also extends.
The positive-current and character equations proving closure
are unchanged.  If $w\ne x$, the section $g$ is regular at $w$;
its standard central splitting acts trivially on the vacuum
generator there.  The Fock vacuum and determinant formula use
the same splitting.  Thus these cycles are exactly the external
highest cycles with the vacuum at $w$, without a factor depending
on $w$.  Their classes generate also when $w=x$, by the point
reduction theorem.

Here is the specialization argument that identifies the actual
comparison map on the diagonal.  Pass to a parameter DVR transverse
to $w=x$, with uniformizer $\pi$, and denote either shifted BRST
complex by $K$.  Its terms are flat; its generic and closed fibers
have cohomology only in degree zero.  The closed degree-zero
classes lift to $K$, because the displayed cycles and their
$W$-modes generate them.  The long exact sequence of
\[
                      K\xrightarrow{\pi}K\longrightarrow K/\pi
\]
therefore shows that multiplication by $\pi$ is injective on
$H^0(K)$ and on every $H^q(K)$.  For $q=1$ use precisely the
surjectivity of the lifting of $H^0(K/\pi)$; in other degrees
injectivity follows directly from the vanishing of the adjacent
fiber group.  The generic calculation makes $H^q(K)$ torsion
for $q\ne0$, hence zero by this injectivity.  The module $H^0(K)$
is torsion-free and consequently flat over the DVR.

On the generic fiber the comparison is the fixed line comparison
at $x$ and the identity on the vacuum.  Since the generating cycles
extend and the degree-zero modules are torsion-free, the same
formula holds before specialization and at $w=x$.  This proves
the comparison on the extremal-label vector.  Degree-zero maps
between the matched label objects are just maps of the
corresponding lines with connection, so it proves the comparison
on the full two-point test locus.  Auxiliary parameters can be
included throughout.  All elementary diagonal comparisons
therefore agree; their repeated composites are coherent by the
closed-push and discrete-mapping realization argument.  A crossing
between the two colors requires no operation identifying their
colors: it is simply parameter pull of the comparison on the
independent marked-point chart.

\smallskip
\noindent\emph{Insertion for the one-slot equivalence.}
Compare $\alpha_V\operatorname{ins}$ with
$\operatorname{ins}\alpha_T$ using \eqref{eq:L11} and
\eqref{eq:L12} with the other color empty.  After applying $F_1$
at $V$, the comparison is the one already constructed.  On heart
objects the gap makes $F_1$ fully faithful in degree zero and
negative degrees, so this lifts uniquely before $F_1$; insertion
is $t$-exact by Lemma~\ref{local:insertion}.  Realization extends
the lifted comparison.  For transitivity, test on the dense
separated locus.  In the test category insertion by a closed
pole image is fully faithful, since $i^!i_*\simeq\operatorname{Id}$
on each closed bound and hence after passage to the colimit.
Thus the same uniqueness using KL generators proves transitivity.
It also proves repetition and the auxiliary linearity coherences.
All unit conventions, determinant lines, and coordinate transports
are retained in the comparisons, as required.
\end{proof}
